\section{Related Work (Jinyu)}

BarrierShare is related to three lines of work, distinguished primarily by their information structure.
First, multi-stage bandits study terminal-only feedback in distributed hierarchies, but use the realized outcome only through local bandit updates.
Second, feedback-graph, semi-bandit, combinatorial bandit, and partial-monitoring models enrich the current-round observation, whereas BarrierShare keeps terminal-only observations and controls the propagation of the induced update.
Third, workflow-constrained agents motivate fixed-stage policies and terminal evaluation, but typically optimize execution within a trajectory rather than repeated online selection among workflow policies.
Table~\ref{tab:feedback_regimes} summarizes these regimes.

\begin{table}[t]
\centering
\small
\caption{Feedback regimes relevant to BarrierShare. Unlike stronger-feedback models, BarrierShare does not enrich the current-round observation; it restricts the scope over which the terminal-cost-induced update trace may be reused.}
\label{tab:feedback_regimes}
\begin{tabular}{lllll}
\toprule
Regime & Observation & Reused object & Reuse scope & Example \\
\midrule
End-to-end bandit & sampled terminal cost & none & none & \citet{hou2024distributed} \\
Complete / one-hop & child/action-level feedback & current-round estimates & local node & Hou complete feedback \\
Full-share endpoint & sampled terminal cost & weight delta & all ancestors & conceptual endpoint \\
BarrierShare & sampled terminal cost & weight delta & safe subtrees & ours \\
\bottomrule
\end{tabular}
\end{table}

\subsection{Multi-Stage Bandits and Distributed Online Learning}

Classical adversarial bandits study a single learner that selects one action and observes only the loss of the selected action, with algorithms such as EXP3 providing no-regret guarantees in the one-stage setting~\citep{auer2002nonstochastic}.
Multi-stage systems introduce a qualitatively different difficulty: the terminal outcome is determined by a sequence of decisions across stages, and upstream decisions affect not only immediate performance but also whether downstream learners receive enough traffic to learn.
This creates the education--exploration--exploitation trilemma formalized by \citet{hou2024distributed}.

The closest theoretical baseline for BarrierShare is Hou's multi-stage bandit model.
In that formulation, a system is represented as a tree: each internal node selects a child, each leaf produces a terminal cost, and the selected path observes only the realized end-to-end outcome.
The $\epsilon$-EXP3 algorithm addresses the education problem by mixing uniform exploration with local exponential-weights updates, yielding regret of order \(O(T^{L/(L+1)})\) for an \(L\)-stage system.
This is the conservative endpoint for our work: terminal feedback is available, but it is used only through local bandit-style updates along the sampled trajectory.

Hou also studies a stronger complete one-hop feedback setting, where each node can observe richer feedback about all of its actions and normalized-EG-type updates become possible.
This setting is fundamentally different from BarrierShare: it enriches current-round observations, whereas our model keeps terminal-only scalar feedback and restricts whether the induced persistent update trace can be shared upward.

The two-stage model of \citet{singla2018learning} is related in that a central forecaster interacts with learning experts and can coordinate their learning by blocking whether the selected expert observes feedback.
However, this feedback blocking controls the learning state of selected experts in a two-level expert-advice problem, whereas BarrierShare studies arbitrary-depth hierarchies in which terminal costs remain locally usable but the induced persistent update delta may propagate only through structurally admissible ancestors.

Our setting is also distinct from adversarial MDP and tree-search formulations.
Although a hierarchy can be viewed as a tree of states, these methods typically assume transition-level state observations, simulator access, value estimates, or centralized coordination.
Unlike adversarial MDP formulations, our learner does not observe intermediate states or rewards along the trajectory, and the regret comparator is a stationary root-to-leaf workflow policy rather than an adaptive Markov policy.
Unlike tree-search methods, our tree is a fixed high-level workflow-policy family rather than a low-level action expansion tree, and learning occurs across repeated terminal-cost observations rather than through within-trajectory search.

\subsection{Feedback Regimes in Hierarchical Learning}

A broad online learning literature studies how regret changes when the learner receives feedback beyond the loss of the selected action.
Full-information learning reveals all action losses; semi-bandit and combinatorial bandit models reveal component-level feedback; partial-monitoring and feedback-graph models specify which side observations become available after each action~\citep{cesa2006prediction, mannor2011bandits, audibert2014regret, alon2015online}.
These settings improve learning by enriching the information available in the current round.

BarrierShare differs in the object being reused.
We do not reveal losses of unselected leaves or child-level costs; the only observed scalar remains the terminal cost of the sampled path.
In a shareable region, this cost induces an exponential-weight delta at the sampled leaf, and this delta updates admissible ancestor aggregates that determine future hierarchical sampling probabilities.
Thus, the question is not which additional losses are observed in the current round, but how far the induced update trace may propagate.

Path and combinatorial bandits also consider actions with internal structure, such as paths, subsets, or other combinatorial objects~\citep{cesa2012combinatorial, audibert2014regret}.
However, semi-bandit variants typically reveal component losses, and bandit variants treat the entire path as a flat arm without providing a mechanism for recursively reusing terminal feedback across safe prefixes.
BarrierShare lies between these extremes: it never observes component-level losses, but it also avoids flattening the hierarchy whenever the induced update is admissible for a safe subtree.

As a useful conceptual endpoint, unrestricted full-share recursive aggregation would allow a selected shareable leaf to update all ancestor aggregates through its weight delta.
BarrierShare retains this mechanism only inside structurally admissible regions; once a barrier is reached, upstream decisions revert to local bandit learning.
Our motivation is also consistent with observations from collaborative and federated learning that shared updates may expose information about data, participants, or model behavior~\citep{melis2019unintended, zhu2019deep, tramer2016stealing, sheller2020federated}.
We do not claim a privacy guarantee: barriers are externally specified sharing constraints, and our goal is to characterize their online-learning consequences.

\subsection{Workflow-Constrained Language Agents}

Recent language-agent systems increasingly structure execution through tools, search, verifiers, memory, or workflow modules.
Tool-using and reasoning-acting agents interleave reasoning, actions, and observations when interacting with external environments~\citep{yao2022react, schick2023toolformer}.
Web-agent and computer-use benchmarks expose long-horizon tasks with tool calls, state tracking, policy constraints, and terminal success criteria~\citep{zhou2023webarena, yao2024tau}.
These works motivate the view that many agent executions are organized around constrained decision points rather than a single flat action.

Importantly, the tree in our formulation is not a low-level action tree over clicks, tokens, or simulator states; it is a fixed high-level policy family whose leaves are complete workflow policies or macro-strategies.
A workflow policy may choose among prompts, tools, specialist modules, API calls, or standard operating procedures at different stages.
Thus, our abstraction is closer to repeated selection among candidate workflow policies than to online search over low-level environment actions.

Recent agent work improves trajectory execution through workflow guidance, process reward models and verifiers, value-guided search, memory retrieval, distillation, or computer-use data engines~\citep{liu2018workflow, chae2025webshepherd, mendes2025selftaught, kang2025agentdistillation, zhang2025gmemory, wang2025opencua}.
These methods change how an agent executes, evaluates, retrieves, or learns within a trajectory.
Unlike prior agent works that improve the executor, verifier, reward model, search procedure, or memory within a trajectory, BarrierShare studies the learning layer above execution: how to repeatedly select among fixed-stage workflow policies when only terminal costs are observed and only part of the induced update trace can be propagated upward.
